%!TEX root = ../main.tex
\begin{titlepage}

\begin{center}
{\LARGE \bf 太陽光発電による電力グリッドに対する\\数理モデルを用いた運用最適化手法}
\\[0.7cm]
{\Large 奥村~侑真}
\\[1.0cm]
{\LARGE \bf 要\vspace{36pt}   旨\\}
\end{center}

近年，地球温暖化対策として「2050年カーボンニュートラル」や「GX2040ビジョン」が掲げられ，再生可能エネルギーの導入が加速している．その中でも太陽光発電は分散型電源として期待されているが，気象条件による出力変動が電力系統の需給バランスに影響を与えるため，蓄電池を活用した電力運用が急務となっている．
特に，市場連動型料金プランの普及に伴い，電力市場の需給に応じた経済的な充放電計画を実運用における時間的制約条件下で策定することが肝要である．本研究では，徳島県内のある工場をモデルとした，太陽光発電設備，蓄電池，および自己託送先を含む電力グリッドシステムを対象とし，その運用コスト削減と効率化を目的とした最適な運用手法を提案する．具体的には，線形計画モデルと動的計画モデルの2つの数理モデルを構築し，それぞれの計算特性や実用性を比較・検証した．線形計画モデルでは，買電コストの最小化や余剰電力の売電益を考慮した正味支出の最小化など複数の目的関数を設定し，厳密解の導出を行った．一方，動的計画モデルでは，蓄電池容量を離散化し，逆方向の最適性の原理に基づくBellman方程式を用いて定式化した．
2025年7月の陽光発電および需要電力の実データを用いた実験の結果，両モデルとも合理的な運用計画を導出したが，計算特性には差異が確認された．線形計画モデルは制約条件の影響を受けやすい一方，動的計画モデルは計画期間に対し計算時間が線形に増加する安定したスケーラビリティを示した．結論として，短期間の高精度な運用には線形計画モデル，長期間の計画や計算時間の予測性を重視する場合には動的計画モデルが有効であるという，実環境における選定指針を明らかにした．
\end{titlepage}
